\section{How a guest's address-map edit becomes real GPU work}
\label{sec:primer}

A GPU virtual address means something only in a particular address space. Its page tables
translate that address to backing memory and permissions. Kayfabe keeps the virtual names the
guest expects, while the host driver builds the real GPU's tables. A \emph{host twin} is the
host-side address space or channel created to correspond to a guest object; it is not a copy
of the GPU.\src{architecture}\src{hostcode}

\subsection{The update path, end to end}

\doc The guest writes page-table entries into GPU memory or guest RAM. Those are ordinary
memory stores through untrapped apertures: Kayfabe does not receive a callback for each edit.
The guest later requests that translations be made current, through an invalidate register
write or the appropriate memory-operation/control path. A register write into BAR0, the
register window, traps into the VMM; command methods are handled by channel workers instead.
The trap records demand and returns. It does not perform the mapping work.\src{architecture}
\src{translated}\src{memcode}

The \emph{walk} is a traversal of the guest's page tables to discover their current meanings.
Kayfabe uses a GPU kernel to read tables already in GPU memory rather than repeatedly pulling
them through a CPU view. It compares the discovered mappings with host placements previously
confirmed successful. \emph{Apply} means asking the host driver to map and unmap those changed
ranges, then invalidate its translation cache. A \emph{reservation} is a host address range
whose ownership and geometry constrain where those mappings can be installed.
\src{architecture}\src{diffcode}\src{applycode}

An acknowledgement says the required mapping transition is settled. Dependent GPU commands
must not use an old translation merely because an emulator returned early. A
\emph{mapping barrier} is that ordering obligation; a \emph{split} divides a command stream
so later work waits for the required refresh. An \emph{admission gate} withholds new software
reads or submissions while their mapping dependency is unsettled. These terms describe
software ordering, not a global pause of all GPU execution.\src{translated}\src{latency}

\begin{figure}[H]
\centering
\begin{tikzpicture}[font=\small,
  box/.style={draw,rounded corners,align=left,text width=65mm,inner sep=3mm},
  flow/.style={-{Latex[length=2mm]},thick}]
\node[box,fill=blue!5] (guest) at (0,0)
  {\textbf{Guest driver}\\1. Write page tables in guest RAM / video memory.\\
   Ordinary memory writes: no per-entry trap.};
\node[box,fill=blue!5] (barrier) at (0,-2.8)
  {2. Request current translations.\\BAR0 register write traps; command/control paths
   reach their workers.};
\node[box,fill=green!5] (walk) at (8,-2.8)
  {\textbf{VA-manager thread + GPU walk}\\3. Re-read live guest tables.\\
   Diff against confirmed host placements.};
\node[box,fill=green!5] (host) at (8,-5.6)
  {\textbf{Host driver / real GPU}\\4. Apply changed host maps and unmaps.\\
   Invalidate; acknowledge actual results.};
\node[box,fill=blue!5] (dependent) at (0,-5.6)
  {5. Release the dependent operation.\\Old backing retires only when the required
   host acknowledgement permits it.};
\draw[flow] (guest) -- (barrier);
\draw[flow] (barrier) -- (walk);
\draw[flow] (walk) -- (host);
\draw[flow] (host) -- (dependent);
\draw[dashed] (-3.65,-1.3) -- (11.65,-1.3);
\node[anchor=west,font=\scriptsize] at (0.2,-1.05)
  {Table bytes alone do not tell the host that a refresh is complete.};
\end{tikzpicture}
\caption{The missing update path. Acknowledging step 2 before step 4 needs another way to
  prevent stale access; it is not justified by recording the request.\src{architecture}
  \src{translated}\src{latency}}
\label{fig:mapping-path}
\end{figure}


\doc A refused map can leave an address absent and therefore faulting. A refused unmap can
leave old host backing reachable, which is a different risk. The project's rules permit
acknowledging the former absence but do not generally permit pretending the latter retirement
succeeded. The early-acknowledgement alternative is therefore a real design question: it would
need fault/replay or another way to prevent stale access.\src{rulings}\src{applycode}

\subsection{Who manages the resources?}

\doc RM means \emph{resource manager}. The driver allocates objects for memory, address spaces
and channels, then names them with handles. A handle is an object name, not a physical address
or an access right by itself. On the guest side, the CPU driver speaks to firmware on the GPU
System Processor, or GSP. Kayfabe plays that firmware-facing peer, including boot, message
queues and replies, then independently asks the real host resource manager for resources.
It does not run NVIDIA firmware inside the emulated device.\src{architecture}\src{gspcode}

A page-directory base is the root used to find an address space's tables. The space can outlive
that particular root; a recycled root address must not silently identify a new object as the
old one. \emph{Mirror} means constructing host mappings from the guest's declared mappings,
not handing hostile guest tables directly to the real GPU.\src{architecture}\src{memcode}

\subsection{What a channel submits}

A \emph{channel} is an ordered execution context for an engine. A copy engine moves or fills
memory; graphics, compute and video engines perform other work. A timeslice group (TSG) groups
related channels for scheduling. Preserving each stream's bytes does not by itself preserve
its group or shared context.\src{kinds}\src{tsgcode}

A channel's \emph{GPFIFO} is a ring of references to command buffers, called pushbuffers.
A descriptor gives an address and length; the pointed-to buffer contains method packets and
arguments for an engine object. Packets can span descriptors. A software relay that inspects
commands must preserve packet state across those boundaries. A \emph{GP reference} is one
such command-buffer descriptor, not the bytes of the commands.\src{relaydesign}\src{turncode}

Each channel also has shared producer/consumer state, called USERD. The driver appends ring
entries, publishes its write index, GP\_PUT, and rings a \emph{doorbell}: a register notification
to inspect work. The GPU fetch unit advances GP\_GET as it reads entries. GET therefore means
``fetched'', not ``finished''. Commands behind a dependency may remain unexecuted after their
reference was fetched. A \emph{token} selects a channel route when a doorbell is rung; the
write does not prove which guest process sent it.\src{eventloop}\src{trapcode}

\subsection{Dependencies, completion and Windows recovery}

A \emph{semaphore} is a memory value used for dependency ordering: one operation releases a
value and another waits for it. A processed marker witnesses that execution crossed a chosen
point. A \emph{non-stall interrupt} invites software to inspect completion state; it is not
itself a per-packet completion receipt. Host notifications can be engine-wide.
\src{eventloop}\src{kinds}\src{faq}

\doc WDDM, the Windows Display Driver Model, schedules GPU tasks and handles recovery.
Microsoft documents a default two-second preemption wait. On a detected timeout Windows
attempts recovery; a recovered display can report that the driver stopped responding, while
applications may need to recreate their device. Engine and adapter recovery have different
scope. A reported 2.43-second host operation is not enough to prove which timeout occurred:
the scheduler's request time and the affected task must be identified.\src{mskeys}
\src{msrecovery}\src{walknarrow}

An \emph{Xid} is a host NVIDIA error record; Windows adapter Code 43 is an adapter status;
events 153 and 4101 are guest log observations used by the test records. They are different
observation domains, not names for one causal defect. A test must correlate them by time and
lifetime rather than treating one quiet counter as proof of health.\src{winmatrix}\src{faq}

\subsection{Sparse memory is a third state}

\doc A backed mapping points to real storage. A plain-invalid address faults. A sparse
mapping intentionally has no backing while absorbing writes and returning zeros on reads.
The selective sparse-mapping experiment tries to preserve that distinction; turning every
unmapped address sparse would also swallow accesses that should fault. Its finite reservation
budget and fallback can therefore affect compatibility as well as resource cost.
\src{sparse}\src{budget}

\doc UVM means unified virtual memory. Some managed-memory accesses need a page supplied on
first touch and the access replayed. Mirroring mappings already present does not deliver that
fault/replay service. This matters both for the remaining application failures and for the
proposal to acknowledge a mapping barrier early.\src{uvmdesign}\src{status}

\subsection{Three kinds of execution, three meanings of permission}

\doc Kayfabe calls ordinary guest-user execution \emph{Passthrough channels}; this is not VFIO
assignment of the whole GPU. Their host twins execute commands under GPU USER privilege in
restricted host address spaces. \emph{Translated channels} serve guest-kernel work by rewriting
addresses into an emulator-owned privileged space. \emph{Emulated channels} perform work that
belongs entirely to the virtual device. Only the last kind completes wholly in software.
\src{kinds}\src{architecture}

GPU USER privilege limits commands on hardware. The host driver's non-privileged request
interface limits which host operations Kayfabe may ask for. A Unix UID controls host process
credentials. A source or benchmark that establishes one boundary does not establish the other
two, and none is a substitute for fault-containment tests.\src{security}\src{physevidence}
\begin{figure}[H]
\centering
\begin{tikzpicture}[font=\small,
  box/.style={draw,align=center,text width=55mm,minimum height=12mm,inner sep=3mm},
  flow/.style={-{Latex[length=2mm]},thick}]
\node[box,fill=red!6] (guest) at (0,0)
  {\textbf{Untrusted guest}\\User programs AND guest root\\own commands, tables, doorbell writes};
\node[box,fill=green!5] (vmm) at (0,-2.5)
  {\textbf{VMM + Kayfabe}\\Validate ownership and bounds.\\Construct host requests.};
\node[box,fill=green!5] (driver) at (0,-5)
  {\textbf{Trusted host driver}\\Non-privileged client interface\\builds real GPU page tables};
\node[box,fill=blue!5] (gpu) at (7.4,-5)
  {\textbf{GPU hardware}\\USER channels + restricted address spaces\\
   privileged rewritten work separate};
\node[box,fill=yellow!8] (host) at (7.4,0)
  {\textbf{Host desktop / other clients}\\Share the physical GPU.\\
   Fault and fairness effects unmeasured.};
\draw[flow] (guest) -- (vmm);
\draw[flow] (vmm) -- (driver);
\draw[flow] (driver) -- (gpu);
\draw[flow,dashed] (host) -- node[fill=white,align=center,text width=44mm]
  {shared execution / recovery} (gpu);
\draw[dashed,thick] (-3.8,-1.2) -- (10.8,-1.2);
\node[anchor=west,font=\scriptsize] at (-3.8,-1.55) {Guest-to-host trust boundary};
\end{tikzpicture}
\caption{Three permission systems: host request permissions, GPU channel privilege and Unix
  credentials. Host address isolation alone does not prove host-desktop fault containment.
  \src{security}\src{rulings}}
\label{fig:trust}
\end{figure}

